%%%%%%%%%%%%%%%%%%%%%%%%%%%%%%%%%%%%%%%%%%%
\section{Tame fundamental groups of log schemes}
\label{s:fg-pi1-log}
%%%%%%%%%%%%%%%%%%%%%%%%%%%%%%%%%%%%%%%%%%%                                
In this section, we deal with tame fundamental groups of log schemes. For our application to tame fundamental groups of rigid spaces in \cref{s:fg-pi1-rig}, we need to consider certain log schemes which are not necessarily fs, but whose log structure is in a certain sense finitely presented relative to a base divisible valuative monoid. We introduced them in \cite{AHLS-Log} and called them \emph{log schemes of type $\typeVd$}. In the same paper, we extended the theory of the Kummer \'etale fundamental group to such log schemes. The relevant results are reviewed in \cref{ss:log-review}. 

The goal of this section is to introduce the \emph{tame} Kummer \'etale fundamental group $\pitame(X/S, \bar x)$ of a connected log scheme $X$ of type $\typeVd$ relative to a base scheme $S$. It is defined using Kummer \'etale covers which become prime-to-$p$ after pull-back to every strictly henselian valued field with a centre on $S$ (see \cref{def:tame-log}). The main result \cref{thm:pi1tame log is fg} states that if $X$ is of finite type over an algebraically closed field $k$, then $\pitame(X/k, \bar x)$ is finitely generated.

This finiteness result will be obtained in three steps. The first step, the case of trivial log structure, was proved in \cref{thm:naked pi1tame fg}. We next treat the case where the log structure is locally constant by means of the map to the underlying scheme endowed with trivial log structure (see \cref{cor:pi1tame locally constant log is fg}). As is well known, the intuition for this `forget log map' is that of a fibration in tori, and this works well here even for log schemes of type $\typeVd$. The general case is then obtained by formal gluing along the log stratification.

In the final \S\ref{ss:surgery} we explain how the tame fundamental group of certain smooth log schemes over a log point can be computed by a ``surgical procedure'' along the strict normalisation map. This will enable us to show that it is finitely presented in certain cases (\cref{cor:surgery-fp}) and to compute some examples.

%%%%%%%%%%%%%%%%%%%%%%%%%%%%%%%%%%%%%%%%%%%%%%%%
\subsection{Review of logarithmic geometry beyond fs} 
\label{ss:log-review}

As we explained in the introduction, dealing with semistable reduction and tamely ramified covers over general valuation rings necessitates the extension of logarithmic geometry beyond the case of fs log schemes. This has been handled in the companion paper \cite{AHLS-Log}, where in particular it was observed that for the most natural definition of Kummer \'etale maps, the underlying maps of schemes might not be of finite type in general. Fortunately for us, this difficulty does not occur if one considers log schemes of type $\typeVd$ as defined in \cite[\S 5]{AHLS-Log} and recalled below. In particular, the theorems and definitions reviewed here are not exactly the general ones stated in \cite{AHLS-Log}, but their simplified versions in the case of log schemes of type $\typeVd$.

Following general conventions (see \cite{Kato1989:LogarithmicStructures,Ogus,KatoToricSingularities} in addition to \cite{AHLS-Log}), the underlying scheme of a log scheme $X$ is denoted by $\underline{X}$, and $\overline{\cM}_{X}$ is the quotient $\cM_X/\cO_X^\times$. We will only consider log schemes which \'etale locally admit a chart by a saturated (but not necessarily finitely generated) monoid.  

Recall that a monoid $V$ is called {\bf valuative} if it is integral and if $V^\gp = V \cup (-V)$. We call $V$ {\bf divisible} if the associated group $V^\gp$ is divisible. 

\begin{defi} \label{def:log-sch-type-Vd}
    A saturated monoid is of {\bf type $\typeVd$} if it is finitely generated over a divisible valuative submonoid. A saturated log scheme is of {\bf type $\typeVd$} if \'etale locally it admits a chart by a monoid of type $\typeVd$.
\end{defi}

Trivially, every fs monoid is of type $\typeVd$, and hence every fs log scheme is of type $\typeVd$. As we shall see, another natural source of such log schemes are semistable schemes over a valuation ring $K^+$ with algebraically closed field of fractions, endowed with the standard log structure (see \cref{ss:log-formal-schemes}). In this case, the base divisible valuative monoid is the non-negative part $\Gamma_K^+$ of the value group of $K$.

\begin{defi}[{\cite[Definition~2.4.4 and Lemma~2.6.1]{AHLS-Log}}] \label{def:Kummer-etale-monoid}
    Fix a set of primes $\Sigma$. Let $P$ be a monoid of type $\typeVd$ and let $\vtheta\colon P\to Q$ be a homomorphism to a saturated monoid $Q$. We say that $\vtheta$ is {\bf Kummer \'etale} if it is injective and exact, and if the quotient $Q^\gp/P^\gp$ is a finite group of order prime to all $p\in \Sigma$. We say that $\vtheta \colon P\to Q$ is {\bf smooth} (resp.\ {\bf \'etale}) if $Q$ is finitely generated over $P$ and the kernel and the torsion part of the cokernel (resp.\ the kernel and the cokernel) of $P^{\rm gp}\to Q^{\rm gp}$ are finite groups of order prime to all $p\in \Sigma$.
\end{defi}

We showed \cite[Lemma~2.6.2]{AHLS-Log} that if $P$ is of type $\typeVd$ and if $P\to Q$ is Kummer \'etale, then $Q$ is finitely generated as a $P$-set, and in particular finitely generated (in fact, even finitely presented) as a monoid over $P$. Consequently, $Q$ is of type $\typeVd$ as well, and the morphism of schemes $\Spec(\bZ[Q])\to \Spec(\bZ[P])$ is finite and of finite presentation. 

\begin{defi} \label{def:Kummer-etale-map}
    Let $X$ be a log scheme of type $\typeVd$ and let $f\colon Y\to X$ be a morphism of saturated log schemes. 
    \begin{enumerate}[(a)]
        \item 
        We say that $f$ is {\bf Kummer \'etale} if \'etale locally on $X$ and $Y$ it admits a chart by a Kummer \'etale morphism of saturated monoids $P\to Q$ for $\Sigma$ the set of primes non-invertible on $X$ such that the induced morphism of schemes
        \[ 
            \underline{Y} \la \underline{X}\times_{\Spec(\bZ[P])} \Spec(\bZ[Q])
        \]
        is \'etale (see \cite[Definition~3.5.1]{AHLS-Log}).
        \item 
        We say that $f$ is {\bf finite Kummer \'etale} if it is Kummer \'etale and if its underlying morphism of schemes is finite (see \cite[Definition~4.2.1 and Corollary~5.1.2]{AHLS-Log}).
    \end{enumerate}
\end{defi}

We note that if $X$ is a log scheme of type $\typeVd$ and if $Y\to X$ is Kummer \'etale, then $Y$ is of type $\typeVd$ as well, and the underlying morphism of schemes $\underline{Y}\to \underline{X}$ is locally quasi-finite and locally of finite presentation \cite[Proposition~5.1.2]{AHLS-Log}. Consequently, if $\underline{X}$ is locally noetherian, then so is $\underline{Y}$. In particular, if moreover $Y$ is quasi-compact, then $Y$ has a finite number of connected components.

A {\bf log geometric point} of a log scheme $X$ is a morphism of log schemes $\bar{x}\to X$ where $\bar{x} = \Spec(P\to k)$ is the spectrum of a separably closed field $k$ endowed with a log structure charted by a monoid $P$ such that $P^\gp$ is $n$-divisible for every $n\geq 1$ invertible in $k$. Every non-empty Kummer \'etale log scheme over such an $\bar{x}$ has a section.

\begin{thm}[{\cite[Theorem~4.3.2]{AHLS-Log}}] \label{thm:FKEt-Galois}
    Let $X$ be a connected log scheme of type $\typeVd$ and let $\bar{x}\to X$ be a log geometric point. Then the category $\FEt_X$ of finite Kummer \'etale morphisms $Y\to X$ is a Galois category on which $\bar{x}$ induces a fibre functor $F_{\bar{x}}$. 
\end{thm}

\begin{defi}[{\cite[Definition~4.3.3]{AHLS-Log}}] \label{def:FKEt-pi1}
    Let $X$ be a connected log scheme of type $\typeVd$ and let $\bar{x}\to X$ be a log geometric point. The fundamental group of the Galois category $\FEt_X$ endowed with the fibre functor $F_{\bar{x}}$ is called the {\bf Kummer \'etale fundamental group} of $X$, and is denoted by $\pi_1(X, \bar{x})$.
\end{defi}

Finite Kummer \'etale maps can be defined locally on the base. More precisely, for a log scheme $X$ of type $\typeVd$, the association $U\mapsto \FEt_U$ forms a stack in the (strict) \'etale topology, and even in the Kummer \'etale topology \cite[Corollary~4.3.7]{AHLS-Log}. 
Another useful fact is its ``topological invariance''.

\begin{prop}[{\cite[Proposition~4.2.4(d)]{AHLS-Log}}] \label{prop:FKEt-top-inv}
    Let $i\colon X_0\to X$ be a strict morphism between log schemes of type $\typeVd$ whose underlying morphism of schemes is a universal homeomorphism. Then the induced functor
    \[ 
        i^*\colon \FEt_X \la \FEt_{X_0}
    \]
    is an equivalence. Consequently, if $X$ is connected and $\bar x\to X$ is a log geometric point, we have an isomorphism
    \[ 
        i_* \colon \pi_1(X_0, \bar{x}) \isomto \pi_1(X, \bar{x}).
    \]
\end{prop}

\begin{defi}
\label{defi:pi1-of-monoid}
    For a (sometimes implicit) choice of a separably closed field $k$, we define $\widehat{\bZ}'(1) = \varprojlim_n \mu_n(k)$ to be the Tate module of $k^\times$.
    For a monoid $P$ we set
    \[
        \pi_1(P) = \Hom(P^{\rm gp}, \widehat{\bZ}'(1)).
    \]
\end{defi}

The notation $\pi_1(P)$ is motivated by the following result.

\begin{prop} \label{prop:FKEt-pi1-local}
    Let $X = \Spec(P\overset{\alpha}{\to} A)$ where $P$ is a monoid of type $\typeVd$ and $A$ is a strictly henselian local ring with residue field $k$. Let $x$ be the closed point of $\underline{X}$, and let $\bar{x}$ be a log geometric point above $x$. In this case, we have $\overline{\cM}_{X,x} = P/\alpha^{-1}(A^\times)$.  Then there is a natural isomorphism of profinite groups
    \[ 
        \pi_1(X, \bar{x}) \simeq \pi_1(\overline{\cM}_{X,x}).
    \]
    In particular, $\pi_1(X, \bar{x})$ is a free $\hZ'$-module of finite rank, so abelian, torsion-free, of order coprime 
    to the residue characteristic at $\bar{x}$, and finitely presented as a profinite group.
\end{prop}
\begin{proof}
All but the group theoretic properties of $\pi_1(\overline{\cM}_{X,x})$ have been proven in \cite[Proposition~4.3.1]{AHLS-Log}.
Since the monoid $\overline{\cM}_{X,x}$ is of type $\typeVd$ and sharp, the group 
$\overline{\cM}_{X,x}^{\rm gp}$ is an extension of a divisible group and a finitely generated free abelian group $A$. Since a divisible group does not admit non-zero maps to $\widehat{\bZ}'(1)$, we deduce that
$\pi_1(\overline{\cM}_{X,x})$ is isomorphic to $\Hom(A,\hZ'(1))$, a free $\widehat{\bZ}'$-module of finite rank.

That $\pi_1(\overline{\cM}_{X,x})$ is finitely presented as a profinite group follows from 
\cref{lem:abelian-fg-implies-finitely-presented} below.
\end{proof}

\begin{lem} \label{lem:abelian-fg-implies-finitely-presented}
    Every finitely generated abelian profinite group $A$ is finitely presented as a (potentially nonabelian) profinite group.
\end{lem}

\begin{proof}
    Let $\psi \colon \hZ^n \surj A$ be a surjection with kernel $B$. Since for every prime $\ell$, pro-$\ell$ completion $(-)^{\wedge \ell}$ is exact on abelian profinite groups, the group $B^{\wedge \ell}$ is the kernel of $\bZ_\ell^n \surj A^{\wedge \ell}$. Since $\bZ_\ell$ is a principal ideal domain, the $\bZ_\ell$-module $B^{\wedge \ell}$ is free of rank at most $n$ and can thus be generated by $n$ elements. We fix a surjection $\bZ_\ell^n \surj B^{\wedge \ell}$ for every $\ell$. Combining these in the product over all $\ell$ we obtain a surjection $\hZ^n \surj B$, and so a finite presentation $\hZ^n \to \hZ^n \to A \to 0$ as $\hZ$-module. 

    To obtain a presentation as a profinite group we take generators $x_i$ for $i=1,\ldots,n$ that we map to the images $\psi(e_i) \in A$ of the standard generators $e_i \in \hZ^n$. As relations among the $x_i$ it suffices to take the pairwise commutators $[x_i,x_j]$ and profinite words representing the finitely many generators of $B \subseteq \hZ^n$. 
\end{proof}

We give two further results on the structure of finite Kummer \'etale maps.

\begin{lem}[{See \cite[Lemma~4.2.7]{AHLS-Log}}] 
\label{lem:fKet-local-structure}
    Let $Y\to X$ be a finite Kummer \'etale map and let $P\to \cM(X)$ be a chart by a~saturated monoid $P$. Then \'etale locally on $X$ there exists a finite collection of Kummer \'etale maps $P\to Q_j$, for $j=1,\ldots, r$, and an isomorphism over $X$ 
    \[
        Y \simeq \coprod_{j=1}^r X\times_{\Spec(P\to \bZ[P])}\Spec(Q_j\to\bZ[Q_j]).
    \]
\end{lem}

\begin{cor} \label{cor:fKet-local-structure}
    In the situation of \cref{lem:fKet-local-structure}, if $X$ is quasi-compact, then there exists a Kummer \'etale homomorphism of monoids $P\to Q$ such that in the (saturated) base change diagram
    \[ 
     \begin{tikzcd}
         Y_Q \ar[r] \ar[d] & X_Q \ar[d] \ar[r] & \Spec(Q \to \bZ[Q]) \ar[d] \\
         Y \ar[r] & X \ar[r] & \Spec(P \to \bZ[P])
     \end{tikzcd}
    \]
    the map $Y_Q\to X_Q$ is strict \'etale. 
\end{cor}

\begin{proof}
By \cref{lem:fKet-local-structure} and the quasi-compactness of $X$ there exists an \'etale covering $\{X_i\to X\}_{i=1}^n$ such that the assertion of said lemma holds over $X_i$. Let us denote the resulting monoids by $Q_{ij}$, for $j=1, \ldots, r(i)$. To obtain the assertion of the corollary, it suffices to take $P\to Q$ such that each of these $P\to Q_{ij}$ factors as $P\to Q_{ij}\to Q$. This is clearly possible, as we can take $Q$ to be the saturation of $P$ in the pushout of $\bigoplus P^{\rm gp}\to \bigoplus Q_{ij}^{\rm gp}$ along the sum map $\bigoplus P^{\rm gp}\to P^{\rm gp}$ (which is an extension of $\bigoplus (Q_{ij}^{\rm gp}/P^{\rm gp})$ by $P^{\rm gp}$).  
\end{proof}


%%%%%%%%%%%%%%%%%%%%%%%%%%%%%%%%%%%%%%%%%%%%%%%%
\subsection{The tame fundamental group of a log scheme}

We now introduce a notion of tameness for finite Kummer \'etale maps of log schemes $Y\to X$ relative to a base scheme $S$, which lets us define the tame fundamental group $\pitame(X/S,\ov{x})$ if $X$ is connected. In the case of schemes, tameness ultimately is defined in terms of extensions of strict henselisations being of degree prime-to-$p$, see \cref{defi:tame-extension-of-valued-fields}. We take the same approach in the case of log schemes, replacing the Galois group of a strictly henselian valued field $(K,K^+)$ with the category of finite Kummer \'etale covers of the log scheme $\Spec(K)$ endowed with the log structure pulled back from $X$. 

\begin{defi} \label{def:prime-to-p-obj}
    Let $\cC$ be a Galois category and let $p\geq 1$ be a natural number. We say that an object $Y$ of $\cC$ is {\bf prime-to-$p$} if there exists a finite collection of Galois objects $Y_i$ of degree\footnote{The degree of an object of a Galois category is the cardinality of its image under a fibre functor.} prime to $p$ and a surjection $\coprod Y_i\to Y$. 
\end{defi}

\begin{rmk} \label{rmk:prime-to-p-obj}
    Let $\Gamma$ be a profinite group. The category of prime-to-$p$ objects in $\Gsets{\Gamma}$ coincides with the category of $\Gamma'$-sets where $\Gamma \surj \Gamma'$ is the prime-to-$p$ completion of $\Gamma$, and for an open subgroup $H\subseteq \Gamma$, the object $\Gamma/H$ is prime-to-$p$ if and only if $H$ is open in the prime-to-$p$ topology \cite[\S 3.1]{RibesZalesskii}. It follows that the full subcategory $\cC'\subseteq\cC$ consisting of prime-to-$p$ objects is a Galois subcategory in general. Further, every map of Galois categories preserves prime-to-$p$ objects, and for a Galois subcategory $\cC_0\subseteq \cC$, an object of $\cC_0$ is prime-to-$p$ if and only if it is prime-to-$p$ when considered as an object of $\cC$. Moreover, for a pair of objects $Y,Z\in\cC$ with $Y$ prime-to-$p$ and non-empty (i.e.\ non-initial), the object $Z$ is prime-to-$p$ if and only if the product $Y\times Z$ is prime-to-$p$.
\end{rmk}

\begin{lem} \label{lem:wreath}
    Let $Z\to Y$ be a morphism in a Galois category $\cC$ and let $p\geq 1$ be a natural number. Consider the following conditions.
    \begin{enumerate}[(i)]
        \item $Z$ is prime-to-$p$.
        \item $Y$ is prime-to-$p$.
        \item for every connected component $W$ of $Y$, the base change $Z\times_Y W$ is prime-to-$p$ as an object of the slice category $\cC_{/W}$ (which is a Galois category).
    \end{enumerate}
    Then (ii)\&(iii)$\Rightarrow$(i)$\Rightarrow$(iii), and if $Z\to Y$ is surjective, then also (i)$\Rightarrow$(ii).
\end{lem}

\begin{proof}
We only need to show (ii)\&(iii)$\Rightarrow$(i). For this, we may assume that $Y=W$ and $Z$ are connected. Write $\cC = \Gsets{\Gamma}$ for a profinite group $\Gamma$; we may assume $Y=\Gamma/H$ for a subgroup $H\subseteq \Gamma$ which is open in the prime-to-$p$ topology on $\Gamma$, and $Z = \Gamma/K$ for an open subgroup $K\subseteq H$. In this case, \cite[Lemma~3.1.4]{RibesZalesskii} implies that $K$ is open in the prime-to-$p$ topology on $H$ if and only if it is open in the prime-to-$p$ topology on $\Gamma$.
\end{proof}

\begin{defi} \label{def:tame-log}
    Let $X$ be a saturated log scheme and let $\underline{X}\to S$ be a map of schemes. 
    We say that a finite Kummer \'etale map $Y\to X$ is {\bf tame relative to $S$} if for every commutative square
    \begin{equation} \label{eqn:tame-log-test-sq}
        \begin{tikzcd}
            \Spec(K)\ar[r] \ar[d] & \underline{X}\ar[d] \\
            \Spec(K^+) \ar[r] & S
        \end{tikzcd}
    \end{equation}
    where $(K, K^+)$ is a strictly henselian valued field and, with $T$ as an abbreviation of $\Spec(K)$ with the pull-back log structure from $X$, the pull-back $Y_T \in \FEt_{T}$ is a prime-to-$p$ object where $p$ is the residue characteristic exponent of $K^+$. 
    
    The full subcategory of Kummer \'etale covers $\FEt_X$ consisting of those which are tame relative to $S$ is denoted by $\FEtt_{X/S}$.
\end{defi}

\begin{lem} \label{lem:tame-log-properties}
    Let $X$ be a saturated log scheme and let $\underline{X}\to S$ be a morphism of schemes. 
    \begin{enumerate}[(a)]
        \item \label{lemitem:tame-log-properties-basic} The full subcategory $\FEtt_{X/S}\subseteq \FEt_X$ is stable under fibre product and subquotients. For finite Kummer \'etale maps $Z\to Y\to X$ with $Y\to X$ tame relative to $S$, we have that $Z\to X$ is tame relative to $S$ if and only if $Z\to Y$ is tame relative to $S$.
        \item \label{lemitem:tame-log-properties-Galois} If $X$ is connected, then $\FEtt_{X/S}$ is a Galois subcategory of $\FEt_X$.
        \item \label{lemitem:tame-log-properties-bc} Given a map of saturated log schemes $X'\to X$ and a commutative square 
        \begin{equation} \label{eqn:tame-log-properties-bc-square}
            \begin{tikzcd}
                \underline{X}'\ar[r] \ar[d] & \underline{X} \ar[d] \\
                S' \ar[r] & S,
            \end{tikzcd}
        \end{equation}
        if $Y\to X$ is tame relative to $S$, then $Y'=Y\times_X X'\to X'$ is tame relative to $S'$. 
        \item \label{lemitem:tame-log-properties-no-log-str} If $X$ has trivial log structure, then $Y\to X$ is tame relative to $S$ if and only if it is tame relative to $S$ in the sense of \cref{defi:tame cover of schemes X over S}.
        \item \label{lemitem:tame-log-properties-tame-local} Let $U\to X$ be a strict \'etale map which is a covering for the Nisnevich topology.
        Then $Y\to X$ is tame relative to $S$ if and only if the base change $Y_U\to U$ is tame relative to $S$.
        \item \label{lemitem:tame-log-properties-strict} A strict finite \'etale map $Y\to X$ is tame relative to $S$ if and only if the underlying map of schemes is tame relative to $S$ in the sense of \cref{defi:tame cover of schemes X over S}. 
        \item \label{lemitem:tame-log-properties-XequalsS} If $\underline{X}=S$ then $\FEtt_{X/S} = \FEt_X$.
        \item \label{lemitem:tame-log-properties-XQ} Let $P\to Q$ be a Kummer \'etale map of saturated monoids with respect to the set of primes non-invertible on $S$ and let $X\to \Spec(P\to \bZ[P])$ be a chart (a strict morphism of log schemes). Set
        \[
            X_Q =  X\times_{\Spec(P\to\bZ[P])} \Spec(Q\to\bZ[Q]) .
        \]
        Then $Y\to X$ is tame relative to $S$ if and only if the base change $Y_Q=Y\times_X X_Q\to X_Q$ is. 
    \end{enumerate}
\end{lem}

\begin{proof}
\ref{lemitem:tame-log-properties-basic} The first two assertions follow directly from the analogous statements for the categories of prime-to-$p$ objects of $\FEt_T$ for test objects $T$ as in the definition. 

For the last one, suppose first that $Z\to Y$ is tame relative to $S$. We will show that $Z\to X$ is tame relative to $S$. Consider a test square \eqref{eqn:tame-log-test-sq} and $T=\Spec(K)$ with log structure pulled back from $X$. We must check that $Z_T\in \FEt_T$ is prime-to-$p$ (with $p$ the residue characteristic exponent of $K^+$), which by \cref{lem:wreath} is equivalent to the fact that for every connected component $W$ of $Y_T$ (with the induced log structure), the preimage $Z_W = Z_T\times_{Y_T} W$ is prime-to-$p$ as an object of $\FEt_W$. Let thus $W\subseteq Y_T$ be a connected component and let $T'=W_{\rm red}$ with the induced log structure. By topological invariance \cite[Proposition~4.2.4(d)]{AHLS-Log}, we have $\FEt_W \simeq \FEt_{T'}$, so the second condition is equivalent to $Z_{T'}\in\FEt_{T'}$ being prime-to-$p$. On the other hand, since $Y\to X$ is integral, so is $T'\to W\to Y_T\to T$. Since $T'$ is connected and reduced and $T=\Spec(K)$, we have $T'=\Spec(K')$ for an algebraic extension $K'$ of $K$. Choosing a valuation subring $(K')^+\subseteq K'$ extending $K$, we obtain a test square of the form \eqref{eqn:tame-log-test-sq} but for $Z\to Y$ relative to $S$. As $Z\to Y$ is tame relative to $S$, the object $Z_{T'}\in \FEt_{T'}$ is prime-to-$p$ as desired.

Conversely, suppose $Z\to X$ is tame relative to $S$ and consider a test square 
\[
    \begin{tikzcd}
        \Spec(K)\ar[r] \ar[d] & \underline{Y}\ar[d] \\
        \Spec(K^+) \ar[r] & S.
    \end{tikzcd}
\]
This square induces by composition with $Y\to X$ a test square as in \eqref{eqn:tame-log-test-sq}. Let $T'$ (resp.\ $T$) be $\Spec(K)$ with log structure pulled back from $Y$ (resp.\ $X$); we have a map $T'\to T$ which is the identity on the underlying scheme. Then $Z_{T'}=Z\times_{Y} T'$ is a sub-object of $Z\times_X T' = (Z\times_X T)\times_T T'$. By assumption, $Z\times_X T$ is a prime-to-$p$ object of $\FEt_T$, and hence its pull-back to $\FEt_{T'}$ is prime-to-$p$.

\ref{lemitem:tame-log-properties-Galois} 
Follows from \ref{lemitem:tame-log-properties-basic} and \cref{lem:Galois-subcategory-criterion}.

\ref{lemitem:tame-log-properties-bc} 
Consider a test square 
    \[
        \begin{tikzcd}
            \Spec(K)\ar[r] \ar[d] & \underline{X}'\ar[d] \\
            \Spec(K^+) \ar[r] & S'
        \end{tikzcd}
    \]
and let $T'$ be $\Spec(K)$ with log structure pulled back from $X'$. By stacking with \eqref{eqn:tame-log-properties-bc-square} we obtain a test square for $X/S$, and let $T$ be $\Spec(K)$ with log structure coming from $X$. We have a map $T'\to T$ which is the identity on the underlying scheme. Now, the pull-back $(Y')_{T'}$ can be expressed as $Y_T\times_T T'$. Since $Y\to X$ is tame relative to $S$, the object $Y_T$ of $\FEt_T$ is prime-to-$p$, and hence so is its base change to $T'$.

\ref{lemitem:tame-log-properties-no-log-str} 
Clear from the definitions.

\ref{lemitem:tame-log-properties-tame-local} 
Follows since in this case 
every test square \eqref{eqn:tame-log-test-sq} lifts to $U$.

\ref{lemitem:tame-log-properties-strict} 
In this case $Y\to X$ is the pull-back of the \'etale map $\underline{Y}\to\underline{X}$ under $X\to\underline{X}$. Consider a test square \eqref{eqn:tame-log-test-sq}. We have a Galois subcategory $\FEt_{\underline{T}}\subseteq \FEt_T$, and hence $Y_T = (\underline{Y})_{\underline T}\times_{\underline{T}} T$ is prime-to-$p$ as an object of $\FEt_T$ if and only if $(\underline{Y})_{\underline T}$ is prime-to-$p$ as an object of $\FEt_{\underline{T}}$. The equivalence follows formally from this.  

\ref{lemitem:tame-log-properties-XequalsS} 
Consider a test square \eqref{eqn:tame-log-test-sq} and let $T^+$ be $\Spec(K^+)$ endowed with the log structure pulled back from $X$ (which makes sense since $\underline{X}=S$). We can base change $Y\to X$ to $T^+$ before base changing further to $T$. As $K^+$ is strictly henselian, by \cref{prop:FKEt-pi1-local} we see that $\pitame(T^+)$ is prime to $p$. So the base change of $Y$ to $T^+$, and hence to $T$, is prime to $p$.

\ref{lemitem:tame-log-properties-XQ} We may assume that $X=\Spec(K)$ and $S=\Spec(K^+)$ for a strictly henselian valued field $(K,K^+)$.
Now $X_Q\to X$ is a non-empty prime-to-$p$ object (when connected, it is Galois of degree $[Q^{\rm gp}:P^{\rm gp}]$ which is prime to $p$), so $X_Q\to X$ is tame.
By \ref{lemitem:tame-log-properties-basic} we see that $Y_Q\to X_Q$ is tame relative to $S$ if and only if $Y_Q\to X$ is.  Thus $Y\to X$ is prime-to-$p$ in $\FEt_X$ if and only if $Y_Q = Y \times_X X_Q$ is prime-to-$p$, see  \cref{rmk:prime-to-p-obj}.
\end{proof}


\begin{defi} \label{def:FKEt-pi1-tame}
    Let $X$ be a connected log scheme and let $\underline{X}\to S$ be a morphism of schemes. For a log geometric point $\overline{x}$ of $X$, we denote by $\pitame(X/S, \overline{x})$ the fundamental group of the Galois category $\FEtt_{X/S}$ with respect to the fibre functor at $\ov{x}$, called the {\bf tame fundamental group} of $X$ relative to $S$. 
\end{defi}

By \cref{lem:tame-log-properties}~\ref{lemitem:tame-log-properties-bc}, a commutative square \eqref{eqn:tame-log-properties-bc-square} with connected log schemes $X$ and $X'$ induces a continuous map $\pitame(X'/S',\ov x)\to\pitame(X/S,\ov x)$. This map is surjective if $X'\to X$ is an isomorphism, and $\pitame(X/X,\ov{x})\simeq \pi_1(X,\ov{x})$. If $X$ has trivial log structure, we have $\pitame(X/S,\ov{x})\simeq \pitame(\underline{X}/S,\ov{x})$.


%%%%%%%%%%%%%%%%%%%%%%%%%%%%%%%%%%%%%%%%%%%%%%%%
\subsection{Log stratifications}
\label{ss:log-strat}

For our results about the tame Kummer \'etale fundamental group we will need to work with log stratifications. This needs some extra care as the monoids charting our log structures might not be finitely generated. Fortunately, things work well if the underlying scheme is noetherian (and with a global chart). 

Recall that a {\bf face} of an integral monoid $P$ is a submonoid $F\subseteq P$ such that $x+y\in F$ implies $x, y\in F$. A {\bf prime ideal} of $P$ is a subset $\fp\subseteq P$ such that $P + \fp = \fp$ and $x+y\in \fp$ implies $x\in \fp$ or $y\in \fp$. The quotient $P/F$ is the quotient of the localisation $(P+F^{\rm gp})/F^{\rm gp} \subseteq P^{\rm gp}/F^{\rm gp}$. Faces of $P$ are in an inclusion-reversing bijection $F \mapsto P \setminus F$ with prime ideals of $P$. The set of all prime ideals of $P$ is denoted by $\Spec(P)$, which we endow with the topology generated by the sets
\[
    D(x) = \{\fp \in \Spec(P) \ | \ x \notin \fp\} \qquad (x\in P).
\]
The space $\Spec(P)$ is a $T_0$-space. 

We proved in \cite[Corollary~2.3.3]{AHLS-Log} that if $\vtheta\colon P\to Q$ is an sfp (e.g.\ smooth) morphism of saturated monoids, then the fibres of the map 
\[ 
    \Spec(Q)\la \Spec(P), \quad \fq \mapsto \vtheta^{-1}(\fq)
\]
are finite; if $P\to Q$ is Kummer \'etale, this map is a homeomorphism. In particular, if $P$ is a monoid of type $\typeVd$, i.e.\ finitely generated over a divisible valuative submonoid $V$, 
then the fibres of $\Spec(P)\to \Spec(V)$ are finite. Hence, $P$ has finitely many faces if $V$ does. For any monoid $P$, we have a natural continuous map
\begin{equation} \label{eqn:SpecZP-to-SpecP} 
    \Spec(\bZ[P]) \la \Spec(P), \quad \fp \mapsto \fp\cap P.
\end{equation}

Consider a log scheme $X$ that admits a global chart $\alpha\colon P\to \cM(X)$ where $P$ is a saturated monoid. The chart corresponds to a strict map
\[ 
    X \la \Spec(\bZ[P])
\]
which we compose with \labelcref{eqn:SpecZP-to-SpecP} to obtain a continuous map
\[ 
    \ov\alpha\colon X \la \Spec(P), \quad x \mapsto P \setminus \alpha^{-1}_x(\cO_{X,x}^\times)
\]
where $\alpha_x$ is the composition $P\to \cM(X) \to \cM_{X,x}$. The fibres of $\ov\alpha$ are called the {\bf log strata} of~$X$. They are locally closed if $P$ has finitely many faces. With no such assumption on $P$, we have the following result in the noetherian case.

\begin{lem} \label{lem:noeth-faces}
    Let $X$ be a log scheme admitting a global chart by a monoid $P$ of type $\typeVd$ (or more generally which is sfp over a finitely generated valuative submonoid). Suppose that the underlying scheme of $X$ is noetherian. Then the image of $\ov\alpha\colon X\to \Spec(P)$ is finite, and its fibres are locally closed subsets of $X$. 
\end{lem}

\begin{proof}
Let $V\subseteq P$ be a divisible valuative submonoid such that $P$ is finitely generated over $V$. Since the fibres of the map $\Spec(P)\to \Spec(V)$ are finite, to show the first assertion it suffices to show that the image of $X\to\Spec(V)$ is finite. 
Since $X$ is noetherian, it has finitely many irreducible components, and therefore it suffices to treat the case when $X=\Spec(R)$ is the spectrum of a noetherian domain $R$. We shall see that in this case the image of $X\to \Spec(V)$ consists of at most two points. 

Suppose that the image of $\Spec(R)\to \Spec(V)$ contains three points, corresponding to prime ideals $\fp_0,\fp_1,\fp_2$ of $V$. Since the prime ideals of $V$ are totally ordered, we may assume that $\fp_0\subsetneq \fp_1\subsetneq \fp_2$ and can choose elements $x\in \fp_1\setminus \fp_0$ and $y\in \fp_2\setminus \fp_1$. 

We claim that $x - ny$ belongs to $V$ for all $n\in\bZ$. Indeed, it suffices to show that $-x + ny$ does not belong to $V$. Let $F_i = V\setminus\fp_i$ denote the respective faces, and let $W = V/F_1$, which is a sharp valuative monoid. The quotient map $V\to W$ sends $y$ to zero and $x$ to a non-zero element. This means that the image of $-x+ny$ in $W^{\rm gp}$ does not belong to $W$, and so $-x+ny$ does not belong to $V$. This proves the claim.

Let $\theta\colon V\to R$ denote the monoid map under consideration. Then $\theta(x)$ is infinitely divisible by $\theta(y)$ in $R$ due to  the claim and
\[
\theta(x) = \theta(y)^n \cdot \theta(x-ny).
\]
Let $\fq_i\subseteq R$  be a prime over $\fp_i\subseteq P$, for $i=1,2,3$. Since the element $\theta(y) \in \fq_2$ is not invertible in the local ring $R_{\fq_2}$, by Krull's intersection theorem we obtain $\bigcap_{n \geq 1} \theta(y)^nR_{\fq_2} = 0$ in $R_{\fq_2}$. 
So $\theta(x)$ maps to $0$ in $R_{\fq_2}$. Since $R$ is a domain, we get that $\theta(x) = 0$. This contradicts the fact that $\theta(x)\notin \fq_0$.

For the second assertion, the image $T = \ov\alpha(X)\subseteq \Spec(P)$ is a finite $T_0$-space, and hence for every $y\in T$, the subspace $\{y\}\subseteq T$ is locally closed. Therefore the fibres of $\ov\alpha\colon X\to T$ are locally closed subsets of $X$. 
\end{proof}

%-------------------------------------------------------
\subsection{Locally constant log structure}
\label{ss:loc-const}

We call a log structure $\cM_X$ on a scheme $X$ {\bf locally constant} if the sheaf $\overline{\cM}_X = \cM_X/\cM_X^\times$ is locally constant (in the \'etale topology). Our next goal is to relate the Kummer \'etale fundamental group of a log scheme to that of the underlying scheme in case the log structure is locally constant (see \cref{prop:torus-fibration-pi1} and \cref{cor:pi1tame locally constant log is fg}). This will be a key ingredient in our proof of finite generation in \S\ref{ss:strat}.

\begin{lem} \label{lem:ovalpha--const}
    Let $X$ be a log scheme of type $\typeVd$ and $\alpha\colon P\to \cM(X)$ a chart where $P$ is a monoid of type $\typeVd$, and $\ov\alpha\colon X\to \Spec(P)$ is the induced map. Suppose that $\ov{\alpha}$ maps $X$ to a single point of $\Spec(P)$ corresponding to a face $F\subseteq P$. 
    
    Then the sheaf $\ov{\cM}_X$ is the constant \'etale sheaf associated to the quotient monoid $P/F$. In particular, the log structure on $X$ is locally constant. 
\end{lem}

\begin{proof}
By definition of $\ov\alpha$, for every $x\in X$, the preimage of $\cO_{X,\ov{x}}^\times$ under $\alpha_{\ov x}\colon P\to \cO_{X,\ov{x}}$ is equal to $F$. Thus $\cM_{X,\ov{x}} \simeq P\oplus_F \cO_{X,\ov{x}}^\times$, the map $P\to \cM_X(X)\to \ov{\cM}_X(X)$ factors through $P/F$ and induces a map of sheaves $\underline{P/F}_X\to \ov{\cM}_X$ which is an isomorphism on stalks and hence an isomorphism. 
\end{proof}

Combining the above with the results on log stratifications in the previous section, we obtain the following that we will use in the proof by noetherian induction of \cref{thm:pi1tame log is fg}.

\begin{cor} \label{cor:generically-log-constant}
    Let $X$ be a non-empty log scheme of type $\typeVd$ whose underlying scheme is noetherian. Then there exists a non-empty open subset $U\subseteq X$ on which the log structure is locally constant. 
\end{cor}
\begin{proof}
    The question of finding a Zariski open $U$ on which $\overline{\cM}_X|_U$ is locally constant is \'etale local on $X$. So we may assume that $X$ admits a global chart $\alpha\colon P\to \cM(X)$ as in \cref{lem:ovalpha--const}.
    By \cref{lem:noeth-faces} there is a finite log stratification, and since $\Spec(P)$ is $T_0$, there is a non-empty open stratum $U$. 
    Then the log structure on $U$ is locally constant by \cref{lem:ovalpha--const}.
\end{proof}

We now turn our attention to a log scheme $X$ of type $\typeVd$ with a locally constant log structure.
We moreover assume that $X$ is connected and fix a log geometric point $\tilde{x} \to X$ above a geometric point $\ov{x} \to \underline{X}$. We endow $\ov{x}$ with the induced log structure. We have a homomorphism
\[
     x_\ast \colon \pi_1(\overline{\cM}_{X,\ov x}) = \pi_1(\ov x,\tilde x) \la \pi_1(X, \tilde x).
\]
We think of the morphism forgetting the log structure
\[
    \ep\colon X \la \underline{X}
\]
as a (locally constant) fibration whose fibres are tori with character group $\overline{\cM}_{X,\ov x}^\gp$. Our next goal is to show right exactness of the associated homotopy sequence when the log structure is locally constant. This will rely on the following lemma. 

\begin{lem}
\label{lem:strict-locus-clopen}
    Let $X$ be a log scheme of type $\typeVd$ with a locally constant log structure. Let $Y\to X$ be a Kummer \'etale map. Then the log structure on $Y$ is locally constant and the set 
    \[ 
        \{ y\in Y \,:\, \ov{\cM}_{X,f(\ov y)}\isomto \ov{\cM}_{Y,\ov y} \text{ for some/every geom.\ pt. $\ov y\to \underline{Y}$ above $Y$}\}
    \]
    (called the {\bf strict locus} of $Y\to X$) is both open and closed in $Y$.
\end{lem}

\begin{proof}
The assertion is \'etale local on $X$ and we may replace $Y$ by a direct summand as a Kummer \'etale cover of $X$. We may therefore assume that $X$ admits a global chart $\alpha \colon P \to \cM(X)$ where $P$ is a monoid of type $\typeVd$. By \cref{lem:ovalpha--const} there is a face $F$ of $P$ and an isomorphism $\underline{P/F}_X \isomto \overline{\cM}_X$.

Using the chart lifting property \cite[Proposition~3.5.8]{AHLS-Log}, working locally and replacing $Y$ with a clopen subset we may assume that $Y\simeq X\times_{\Spec(\bZ[P])}\Spec(\bZ[Q])$ for a Kummer \'etale map of monoids $P\to Q$ (with respect to the set of primes non-invertible on $X$). 

We may further assume that $X$ is reduced. Note that in this case the map $\ov\alpha : X \to \Spec(P)$ factors through the subscheme (log stratum of $\Spec(\bZ[P])$)
\[ 
    \Spec(P\to \bZ[F^\gp])
\]
where the map is the inclusion $F\hookrightarrow F^\gp$ and sends everything else to zero (it factors set-theoretically, but since $X$ is reduced, it factors scheme-theoretically). Since for a strict map $g\colon W\to Z$ we have $g^*\ov{\cM}_Z\isomto \ov{\cM}_W$, we may assume that $X = \Spec(P\to \bZ[F^\gp])$, in which case 
\[ 
    Y_{\rm red} = \left(X\times_{\Spec(\bZ[P])}\Spec(\bZ[Q]))\right)_{\rm red} \simeq \Spec(Q\to \bZ[G^\gp])
\]
where $G\subseteq Q$ is the face generated by $F$ (note that since $P\to Q$ is Kummer \'etale, we have $F = P\cap G$). Thus $\overline{\cM}_Y$ is the constant sheaf with value $Q/G$, and $Y$ has locally constant log structure. Moreover, if $P/F\isomto Q/G$, the strict locus is everything, and otherwise it is empty. The strict locus is therefore open and closed.
\end{proof}

We shall combine \cref{lem:strict-locus-clopen} with the following result about the strict locus.

\begin{lem} \label{lem:criterion strict}
    Let $X$ be a log scheme of type $\typeVd$ with a locally constant log structure, let $Y\to X$ be a Kummer \'etale map, and let $\ov x$ be a geometric point of $\underline{X}$.  Then the set of points in the geometric fibre $h^{-1}(\ov x)$ above $\ov x$ at which $h$ is strict agrees with the projection under $h^{-1}(\tilde{x}) \to h^{-1}(\ov x)$ of the 
    fixed points of the natural action of $\pi_1(\ov x,\tilde x)$ on the log geometric fibre $h^{-1}(\tilde x)$.
\end{lem}

\begin{proof}
We may work \'etale locally on $X$ and thus assume that $X$ has a chart by a monoid $P$ of type $\typeVd$.
We may further replace $Y \to X$ by a connected component of its base change to $\ov x = \Spec(P\to k)$. By the structure of Kummer \'etale covers over a strict henselian local ring (\cref{prop:FKEt-pi1-local} or \cite[Lemma~4.2.7]{AHLS-Log}) we may assume that $Y = \Spec(Q\to k\otimes_{\bZ[P]}\bZ[Q])$ for a Kummer \'etale map of monoids $P \to Q$. Let $F \subseteq P$ be the preimage of $k^\times$ and $G\subseteq Q$ the unique face above $P$ (we have $\Spec(Q)\isomto\Spec(P)$ as $P\to Q$ is Kummer \'etale). Then $Y\to X$ is strict if and only if $P/F\isomto Q/G$, which translates into the trivial action of $\pi_1(\ov{x},\tilde{x}) = \pi_1(P/F)$ on $h^{-1}(\tilde x) = Y_{\tilde x}$.
\end{proof}

\begin{prop} \label{prop:torus-fibration-pi1}
    Let $X$ be a connected log scheme of type $\typeVd$ with a locally constant log structure. 
    The map $\ep\colon X \to \underline{X}$ to the underlying scheme induces an exact sequence
    \[
        \begin{tikzcd}
            \pi_1(\ov \cM_{X,x}) \ar[r] & \pi_1(X,\bar x) \ar[r,"\ep_\ast"] & \pi_1(\underline{X},\bar x) \ar[r] & 1.
        \end{tikzcd}
    \]
\end{prop}

\begin{proof}
An \'etale cover $\underline{Y} \to \underline{X}$ induces a strict Kummer \'etale map $Y \to X$. Therefore its base change $Y$ along $\ep$ simply endows $\underline{Y}$ with the pull-back log structure. In particular, the base change of a connected \'etale cover remains connected. This proves that $\ep_\ast$ is surjective. 

To prove exactness at $\pi_1(X, \bar x)$ we must show that a connected \'etale cover $h\colon  Y \to X$ with a section above $\bar x$ (the point being endowed with the pull-back log structure) actually already descends to an \'etale  cover of $\underline{X}$. By \cref{lem:criterion strict}, the existence of the section shows that the locus $S \subseteq Y$ where $Y \to X$ is strict is non-empty. The strict locus $S$ is both open and closed by \cref{lem:strict-locus-clopen}.
Since $Y$ is connected, we must have $S=Y$ and $Y \to X$ is a strict map. Hence $Y \to X$ descends as an \'etale cover $\underline{Y} \to \underline{X}$ and the proof is complete.
\end{proof}

We deduce an analogous result for tame fundamental groups.

\begin{cor} \label{cor:torus-fibration-tame}
    Let $X$ be a connected log scheme of type $\typeVd$ with a locally constant log structure. 
    Let $\underline{X}\to S$ be a morphism of schemes. Then the sequence of tame Kummer \'etale fundamental groups induces an exact sequence
    \[
        \begin{tikzcd}
            \pi_1(\ov \cM_{X,x}) \ar[r] & \pitame(X/S,\bar x) \ar[r,"\ep_\ast"] & \pitame(\underline{X}/S,\bar x) \ar[r] & 1.
        \end{tikzcd}
    \]
\end{cor}

\begin{proof}
Consider the commutative diagram
\[
    \begin{tikzcd}
        \pi_1(\ov \cM_{X,x}) \ar[d,equal] \ar[r] & \pi_1(X,\bar x) \ar[r,"\ep_\ast"]\ar[d,two heads]  & \pi_1(\underline{X},\bar x) \ar[r]\ar[d,two heads] & 1 \\
        \pi_1(\ov \cM_{X,x}) \ar[r,"\iota_{x,\ast}"] & \pitame(X/S,\bar x) \ar[r,"\ep_\ast"] & \pitame(\underline{X}/S,\bar x) \ar[r] & 1.
    \end{tikzcd}
\]
The top row is exact by \cref{prop:torus-fibration-pi1}. The vertical arrows are surjective. Thus the image of $\iota_{x,\ast}$ is normal. Moreover, the map $\varepsilon_*$ in the bottom sequence is surjective. 

To prove that the bottom row is exact, it remains to check that a finite $\pitame(X/S,\ov{x})$-set $M$ on which $\pi_1(\ov \cM_{X,x})$ acts trivially comes from a $\pitame(\underline{X}/S,\bar x)$-set. Translating in terms of covering spaces, let $Y\to X$ be the corresponding finite Kummer \'etale map. By exactness of the top sequence, we see that $Y$ is in the image of $\FEt_{\underline{X}}$, i.e.\ $Y\to X$ is strict. Then \cref{lem:tame-log-properties}~\ref{lemitem:tame-log-properties-strict} implies that $\underline{Y}\to\underline{X}$ is tame relative to $S$, so that $Y\to X$ is in the image of $\FEtt_{\underline{X}/S}$.
\end{proof}

\begin{cor}
\label{cor:pi1tame locally constant log is fg}
    Let $X$ be a connected log scheme of type $\typeVd$ whose underlying scheme is of finite type over an algebraically closed field $k$. Suppose that $X$ has locally constant log structure.
    \begin{enumerate}
        \item 
        The group $\pitame(X/k,\bar x)$ is finitely generated.
        \item 
        The group $\pitame(X/k,\bar x)$ is finitely presented if and only if $\pitame(\underline{X}/k,\bar x)$ is finitely presented.
    \end{enumerate}
\end{cor}

\begin{proof}
    By \cref{prop:FKEt-pi1-local} and 
    \cref{lem:abelian-fg-implies-finitely-presented},
    the image of 
    $\pi_1(\ov \cM_{X,x})$ in $\pitame(X/k,\bar x)$ is a finitely presented profinite group.  
    Thus, by 
    \cite[Proof of Claim~2.7]{EsnaultShustermanSrinivas2022:FinitePresentationTame}, the group 
    $\pitame(X/k,\bar x)$ is finitely generated (resp.\ finitely presented) if and only if $\pitame(\underline{X}/k,\bar x)$ has the same property. 

    In particular, because $\pitame(\underline{X}/k,\bar x)$ is finitely generated by \cref{thm:naked pi1tame fg}, we also know that $\pitame(X/k,\bar x)$ is finitely generated.    
\end{proof}

%%%%%%%%%%%%%%%%%%%%%%%%%%%%%%%%%%%%%%%%%%%%%%%%
\subsection{Finite generation} 
\label{ss:strat}

We shall use \cref{cor:pi1tame locally constant log is fg} (the case of locally constant log structure) as a bootstrap, combined with a formal gluing argument, to prove the following theorem, which is the main result of this section. 

\begin{thm} \label{thm:pi1tame log is fg}
    Let $k$ be an algebraically closed field.
    Let $X$ be a non-empty connected log scheme of type $\typeVd$
    whose underlying scheme is of finite type over $k$.
    Then $\pitame(X/k,\bar x)$ is finitely generated.
\end{thm}

The main difficulty in the proof of \cref{thm:pi1tame log is fg} (deferred to the end of this subsection) hides in the proof of the following proposition, in which a ``formal gluing'' argument is employed. In order to avoid dealing with log rigid analytic spaces, we only formulate it in the affine case.


\begin{prop} \label{prop:pi1tame-strat}
    Let $X=\Spec(R)$ be a non-empty connected affine log scheme of type $\typeVd$ which is of finite type over a field $k$. Let $U=\Spec(R[1/f])$ be a distinguished affine open and let $Z \subseteq X$ be a closed subscheme with $|Z| = |X|\setminus |U|$. We equip both $U$ and $Z$ with the induced log structure. 
    
    Suppose that for every connected component $W$ of either $U$ or $Z$, the tame fundamental group $\pitame(W/k)$ is finitely generated.     
    Then 
    $\pitame(X/k, \bar x)$ 
    is finitely generated. 
\end{prop}
\begin{proof}
Let $T = \Spec(\widehat{R})$ where $\widehat{R} = \varprojlim R/f^{n+1}$ is the $f$-adic completion of $R$. The inclusion $Z \inj T$ induces a bijection $\pi_0(Z) \isomto \pi_0(T)$ of connected components. For $\alpha \in \pi_0(Z)$ corresponding to the connected component $Z_\alpha$  we denote by $T_\alpha$ the connected component of $T$ containing $Z_\alpha$. 

More generally, if $T'\to T$ is a finite morphism, and $Z' = Z\times_T T'$, then $T'$ is the spectrum of a ring which is complete with respect to the ideal $f\cdot \cO_{T'}$, and thus $\pi_0(Z')\isomto\pi_0(T')$ is also a bijection. Since the underlying maps of finite Kummer \'etale covers are finite (as remarked in \cref{def:Kummer-etale-map}, this follows from \cite[Corollary~5.1.2]{AHLS-Log} because $X$ is of type $\typeVd$), this formally implies that, if we endow $T$ with the log structure pulled back from $X$, so that $Z \inj T$ is strict, then for every $\alpha$ and a log geometric point $\overline{z}\to Z_\alpha$, the map
\[
    \pitame(Z_\alpha/k, \overline{z})\la\pitame(T_\alpha/k,\overline{z})
\]
is surjective. In particular, $\pitame(T_\alpha/k,\overline{z})$ is finitely generated. 


Let us also consider $T^\circ = T\times_X U = \Spec(\widehat{R}[1/f])$ that fits into a cartesian square
\[ 
    \begin{tikzcd}
        T^\circ \ar[r,"s"] \ar[d,"c",swap] & U \ar[d] \\
        T \ar[r] & X.
     \end{tikzcd}
\]
Because $U$, $T$ and $T^\circ$ are noetherian, passing to connected components defines a finite graph 
\[
    E_\bullet = \Big[\begin{tikzcd}
        \pi_0(T) \sqcup \pi_0(U) & \ar[l,"c",shift left=0.7ex] \ar[l,"s",swap,shift right=0.7ex] \pi_0(T^\circ) & \ar[l,shift left=1ex] \ar[l,shift right=1ex] \ar[l] \varnothing  
    \end{tikzcd}\Big].
\]

Let us temporarily work with a more general $X$ that is not necessarily connected. Let $\idem(S)$ denote the set of idempotents in $\cO(S)$ for a scheme $S$. By \cite[Theorem~2.2]{ArtinFormalModuliII}, the sequence 
\[ 
    \begin{tikzcd}
        0 \ar[r] &\cO(X) \ar[r] &  \cO(T)\times \cO(U) \ar[r,"-"] &\cO(T^\circ),
    \end{tikzcd}
\]
is exact. Hence, the 
idempotents $\idem(X)$ of $X$ are the equaliser of 
\[
\begin{tikzcd}
        \idem(T) \times \idem(U) \ar[r,"c^\ast",shift left=0.7ex] \ar[r,"s^\ast",swap,shift right=0.7ex] 
        & \idem(T^\circ) .
    \end{tikzcd}
\]
This is equivalent to the map 
\[
\rH^0(X,\bZ/2\bZ) \isomto \rH^0(|E_\bullet|,\bZ/2\bZ)
\]
being an isomorphism. It follows that the natural map $\pi_0(|E_\bullet|) \la \pi_0(X)$
is bijective. In particular, if we again assume that $X$ is connected, then the graph $E_\bullet$ is also connected.


We now define a fibred category in Galois categories $\cG_\bullet$ over $E_\bullet$ by associating 
\begin{itemize}
    \item 
    $\FEtt_{T_\alpha/k}$ to $\alpha \in \pi_0(T)$, and
    \item $\FEtt_{U_\beta/k}$ to $\beta \in \pi_0(U)$ describing the connected component $U_\beta \subseteq U$, and likewise
    \item 
    $\FEtt_{T^\circ_\gamma/k}$ to $\gamma \in \pi_0(T^\circ)$ describing the connected component $T^\circ_\gamma \subseteq T^\circ$. 
    \item 
    The functors $c^*$ and $s^*$ of $\cG_\bullet$ are induced by the maps 
    $T^\circ_\gamma\to T_{c(\gamma)}$ and $T^\circ_\gamma\to U_{s(\gamma)}$. 
\end{itemize}
Because $E_\bullet$ is connected, \cref{prop:DD is galois for connected E} shows that $\Desc(\cG_\bullet)$ is a Galois category. Since $E_\bullet$ is finite and the connected components of $T$ and of $U$ have finitely generated tame fundamental groups, \cref{prop:DD-fgfp} shows that $\pi_1(\Desc(\cG_\bullet))$ is finitely generated. 

Next, we have an obvious pull-back functor between Galois categories
\begin{equation} \label{eqn:tame-desc}
    f^*\colon \FEtt_{X/k} \la \Desc(\cG_\bullet)
\end{equation}
that is exact and thus a morphism of Galois categories. It remains to show that 
the functor \eqref{eqn:tame-desc} is fully faithful. This means that we have to show that for a connected tame cover $X' \to X$, the induced object $f^*(X'\to X)$ in $\Desc(\cG_\bullet)$ is still connected. We use the criterion \cref{lem:local-system-connected} and need to show that with $T' = T \times_X X'$, and $U' = U \times_X X'$ and $T'^\circ = T^\circ \times_X X'$ the graph
\[
E'_\bullet = \Big[\begin{tikzcd}
        \pi_0(T') \sqcup \pi_0(U') & \ar[l,"c",shift left=0.7ex] \ar[l,"s",swap,shift right=0.7ex] \pi_0(T'^\circ) & \ar[l,shift left=1ex] \ar[l,shift right=1ex] \ar[l] \varnothing  
    \end{tikzcd}\Big]
\]
is connected. This follows from the bijection $\pi_0(|E'_\bullet|) \isomto  \pi_0(X')$
we proved above, because $X'$ is connected. This completes the proof. 
\end{proof}

\begin{proof}[Proof of \cref{thm:pi1tame log is fg}]
By descent \cite[Proposition~4.2.4]{AHLS-Log} and by an application of \cref{cor:DD-fgfp}~\ref{coritem:DD-fg}, we may work \'etale locally on $X$ (we have seen this strategy used in the proof of \cref{thm:naked pi1tame fg}). We may therefore assume that $X=\Spec(P \to R)$ is affine with a global chart by a monoid $P$ of type $\typeVd$.

\Cref{cor:generically-log-constant} shows that in the situation of \cref{thm:pi1tame log is fg}, there exists a non-empty open subset $U\subseteq X$ on which the log structure is locally constant. Shrinking $U$, we may assume that $U=\Spec(R[1/f])$ is a distinguished affine open. By \cref{cor:pi1tame locally constant log is fg}, the tame fundamental group of every connected component of $U$ is finitely generated. By noetherian induction, the proof now follows from \cref{prop:pi1tame-strat}.
\end{proof}


%%%%%%%%%%%%%%%%%%%%%%%%%%%%%%%%%%%%%%%%%%%%%%%%
\subsection{Surgery, finite presentation, and examples}
\label{ss:surgery}

In this subsection, we show how the tame fundamental groups of certain log schemes (a class including the special fibres of strictly semistable formal schemes) can be computed by a surgical procedure of decomposing said scheme into its irreducible components and partially stripping off its log structure. This strategy is enabled by the following descent result, generalised beyond the fs setting in \cref{cor:descent-beyond-fs}.

\begin{thm}[{\cite[Theorem~3.2.25]{Stix2002:Thesis}}] \label{thm:Stix-descent}
    Let $f\colon S'\to S$ be an exact morphism of fs log schemes whose underlying map of schemes is finitely presented, proper, and surjective. Then $f$ satisfies universal effective descent for finite Kummer \'etale maps. 
\end{thm}

\begin{cor} \label{cor:descent-beyond-fs}
    Let $f\colon S'\to S$ be a strict morphism of saturated log schemes whose underlying map of schemes is finitely presented, proper, and surjective. Then $f$ satisfies universal effective descent for finite Kummer \'etale maps.
\end{cor}

\begin{proof}
The assertion is (strict) \'etale local on $S$, and hence we may assume that $S$ has a chart by a saturated monoid $P$. Write $P=\varinjlim P_\lambda$ with $P_\lambda$ fs, let $S_\lambda$ be $\underline{S}$ with log structure induced by $P_\lambda$, and let $S'_\lambda$ be $\underline{S}'$ with log structure so that $S'_\lambda\to S_\lambda$ strict. Then by \cite[Proposition~4.2.6]{AHLS-Log} we have
\[
    \FEt_S \simeq \varinjlim \FEt_{S_\lambda}, \qquad \FEt_{S'} \simeq \varinjlim \FEt_{S'_\lambda}
\]
(and the same for $S'\times_S S'$ and $S'\times_S S'\times_S S'$). The maps $S'_\lambda \to S_\lambda$ satisfy the assumptions of \cref{thm:Stix-descent}, and the result follows formally from this. 
\end{proof}

\begin{rmk}
    It is likely that the assertion of \cref{cor:descent-beyond-fs} holds more generally for $f$ exact but not necessarily strict. However, one would need to argue why the maps $S'_\lambda\to S_\lambda$ will be exact, proper, and surjective for $\lambda\gg 0$.
\end{rmk}

We are interested in tame fundamental groups of log schemes which are smooth over a generalised log point. Because of the tameness condition, we have to assume that charts exist Zariski locally, not just \'etale locally. The precise assumptions and notation are listed below.

\begin{setup} \label{setup:surgery}
Consider the following situation: 
\begin{itemize}
    \item 
    $k$ is a field,
    
    \item 
    $S = \Spec(M\to k)$ for a saturated monoid $M$ which is finitely generated over a sharp divisible valuative submonoid $V\subseteq M$ and a map $\alpha\colon M\to k$ with $\alpha^{-1}(k^\times) = \{0\}$ (the key cases of interest are $M=V$ and $V=0$), in particular $M$ is sharp,
    
    \item 
    $X\to S$ is a smooth morphism admitting a chart Zariski locally in the following sense: Zariski locally on $X$ there exists a smooth monoid morphism $\theta\colon M\to P$ (with respect to the set of primes non-invertible in $k$) and a strict \'etale map
    \[
        X \la \Spec(P\to k[P]\otimes_{k[M]} k) = \Spec\big(P\to k[P]/(\theta(M_+))\big)
    \]
    where $M_+ = M\setminus M^\times = M \setminus 0$.
    We assume moreover that $X$ is quasi-compact and connected.
     
    \item $\nu\colon X^{(0)}\to X$ is the normalisation map on the underlying schemes, endowed with the pull-back log structure,
    
    \item $X^{(a)}$ is the $(a+1)$-fold fibre product of $X^{(0)}\to X$, endowed with the pull-back log structure,
    
    \item $U^{(a)}\subseteq X^{(a)}$ is the open subset where $\overline{\cM}_{X^{(a)}}$ (which coincides with the pull-back of $\overline{\cM}_X$ to $X^{(a)}$) is locally constant ($U^{(a)}$ is open and dense by \cref{cor:generically-log-constant}),
    
    \item $X^{(a)}_\star$ is the scheme $\underline{X}{}^{(a)}$ endowed with the log structure induced by the open subset $U^{(a)}$.
\end{itemize}
\end{setup}

\begin{rmk}
As we shall learn in \cref{lem:surgery-properties}, there are natural maps $\varepsilon^{(a)}\colon X^{(a)}\to X^{(a)}_\star$ which intuitively are torus bundles as in \S\ref{ss:loc-const}, and the log schemes $X^{(a)}_\star$ are log smooth over $k$ with no log structure, i.e.\ toroidal embeddings. The following diagram then illustrates how in our setup the log scheme $X$ is disassembled into the log schemes $X_\star^{(a)}$. 
\[
    \begin{tikzcd}
        X \ar[dd,"\text{smooth}",swap] & X^{(0)} \ar[l,swap,"\nu"] \ar[d,"\varepsilon^{(0)}"] & \ar[l,shift right=0.7ex] \ar[l,shift left=0.7ex] X^{(1)} \ar[d,"\varepsilon^{(1)}"] \ar[d] & \ar[l,shift right=1ex] \ar[l,shift left=1ex] \ar[l] X^{(2)} \ar[d] \ar[d,"\varepsilon^{(2)}\text{ ``torus bundle''}"] & \cdots\\
        & X^{(0)}_\star \ar[dr] & X^{(1)}_\star \ar[d] & X^{(2)}_\star \ar[dl,"\text{smooth}"] & \cdots \\
        S = \Spec(M\to k) \ar[rr] & & \Spec(k)
    \end{tikzcd}
\]
\end{rmk}

\begin{ex}[Semistable curve] \label{ex:semistable-log-curve}
    Consider the above setup with $M=V$, let $\pi\in V$ be a nonzero element, and let $P= P_2(\pi) = V[e_1,e_2]/(e_1+e_2=\pi)$ be the ``standard semistable'' monoid as in \cite[Example~2.4.5]{AHLS-Log}. Let $X = \Spec(P\to k[x_1,x_2]/(x_1x_2))$ where the map sends $e_i$ to $x_i$ and all of $V\setminus \{0\}$ to zero. In this case, we have 
    \begin{align*}
        X^{(0)} & = \Spec(P\to k[x_1]) \sqcup \Spec(P\to k[x_2]), \\
        X^{(1)} & = X^{(0)} \sqcup \Spec(P\to k) \sqcup \Spec(P\to k),
    \end{align*}
    and 
    \[
        X^{(0)}_\star  = \Spec(\bN\to k[x_1]) \sqcup \Spec(\bN\to k[x_2]),
        \qquad 
        X^{(1)}_\star = X^{(0)}_\star \sqcup \Spec(k) \sqcup \Spec(k).
    \]
    We notice that $X^{(a)}_\star$ is log smooth over $k$ (with no log structure). This is the case in general.
\end{ex}

\begin{lem} \label{lem:strata-closures}
     In the Zariski local situation of \cref{setup:surgery}, where we consider the log scheme  
    \[
         W = \Spec(P\to k[P]\otimes_{k[M]} k) = \Spec\big(P\to k[P]/(\theta(M_+))\big),
    \]  
    let $F \subseteq P$ be a face in the image of $W \to \Spec(P)$. Then the following hold.
    \begin{enumerate}[(a)]
        \item  \label{lemitem:strata-closures1}
        The monoid $F$ is fs and $0 \to F$ is smooth.
        
        \item \label{lemitem:strata-closures2}
        The log scheme $Z_F = \Spec(P \to k[F])$, where $x \in P \setminus F$ maps to $0$ and $x \in F$ maps to $x \in k[F]$, is the closure of the log stratum of $W$ associated to $F$.
        
        \item  \label{lemitem:strata-closures3}
        The log scheme $Z_{F,\ast} = \Spec(F \to k[F])$ is log smooth over $\Spec(k)$ with trivial log structure, and it agrees with the log scheme structure on $\underline{Z}_F$ induced by the open subset that is the log stratum associated to $F$.
        
        \item  \label{lemitem:strata-closures4}
        The underlying scheme $\underline{Z}_F = \Spec(k[F])$ is normal.
    \end{enumerate}
\end{lem}
\begin{proof}
    \labelcref{lemitem:strata-closures1} The image of $F$ in $\Spec(M)$ by the map $\theta: \Spec(P) \to \Spec(M)$ must be the image of $\Spec(M \to k) \to \Spec(M)$ which is the face $0$. In other words $0 = \theta^{-1}(F) = F \cap M$. The morphism $0 = \theta^{-1}(F) \to F$ is sft by \cite[Lemma~2.3.2]{AHLS-Log}, and hence $F$ is an fs monoid by \cite[Proposition~2.2.9]{AHLS-Log}.    

    Since $M$ is sharp and saturated, $M^\gp$ is torsion free. Hence torsion of $F^\gp$ injects into $P^\gp/M^\gp$ which only contains torsion of order prime to $p$ because $M \to P$ is smooth. It follows that $0 \to F$ is smooth. 

    \labelcref{lemitem:strata-closures2}
    The map $P \to k[F]$ induces a surjective homomorphism $k[P]/(M_+) \surj k[F]$ and thus a strict closed immersion $Z_F \inj W$. The associated vanishing ideal is generated by $P \setminus F$, which is the defining ideal for the closure of the stratum associated to $F$.

    \labelcref{lemitem:strata-closures3} This follows from \labelcref{lemitem:strata-closures1}.

    \labelcref{lemitem:strata-closures4}
    By \cite[Proposition I.3.4.1]{Ogus}, $\underline{Z}_F$ is normal since $F$ is fs. 
\end{proof}

\begin{lem} \label{lem:surgery-properties}
    In the situation of \cref{setup:surgery}, the following statements hold. 
    \begin{enumerate}[(a)]
        \item \label{lemitem:surgery-finite-type}
        The underlying scheme of $X$ is of finite type over $k$.
        
        \item \label{lemitem:surgery-finite-normalisation}
        The map $\nu\colon X^{(0)}\to X$ is finite and surjective, and for every point $x \in X^{(0)}$ the residue fields $k(x) = k(\nu(x))$ agree. 
        
        \item \label{lemitem:surgery-tameness}
        A finite Kummer \'etale map $Y\to X$ is tame relative to $k$ if and only if its base change $Y^{(0)} = Y\times_X X^{(0)}\to X^{(0)}$ is tame relative to $k$. 
                
        \item \label{lemitem:surgery-fs-and-smooth}
        The log scheme $X^{(a)}_\star$ is fs and smooth over $k$ (with trivial log structure), and $U^{(a)}$ is the locus where the log structure on $X^{(a)}_\star$ is trivial.

        \item \label{lemitem:surgery-relative-strip-off-log}
        There exists a unique morphism of log schemes $\varepsilon^{(a)}\colon X^{(a)} \to X^{(a)}_\star$ inducing the identity on the underlying schemes. 
    \end{enumerate}
\end{lem}

\begin{proof}
These assertions are Zariski local on $X$, so we may assume that there exists a smooth morphism of saturated monoids $M\to P$ (necessarily injective) and a strict \'etale map  
\[ 
    X \la W = \Spec(P\to k[P]/(M_+)), \qquad M_+ = M\setminus 0 \subseteq P.
\]
By \cite[Corollary~2.5.18]{AHLS-Log} we see that $P$ is finitely presented over $M$, hence $W$ and consequently $X$ are of finite type, showing \labelcref{lemitem:surgery-finite-type}. In particular, the normalisation map $\nu$ is finite, showing the first part of  \labelcref{lemitem:surgery-finite-normalisation}.

In order to show the claim on the residue fields in \labelcref{lemitem:surgery-finite-normalisation}, we first explicate the map $\nu\colon X^{(0)}\to X$ in the Zariski local situation $W$ used in the previous paragraph and studied in \cref{lem:strata-closures}. 
The irreducible components of~$W$ correspond to the connected components of the $Z_F \inj W$ associated to maximal faces $F$ in the image of $W \to \Spec(P)$. These $Z_F$ are all normal and thus $W^{(0)}$, the normalisation of $W$, is the disjoint union of these maximal $Z_F$. Since $X \to W$ is strict \'etale, the same holds for the normalisation $X^{(0)} = X \times_W W^{(0)}$: this is the disjoint union of the connected components of the $X_F \inj X$ which are the base changes of $Z_F \inj W$ along the strict \'etale $X \to W$; they can also be described as the closures of the open strata of the log stratification of $X$.

Since $X^{(0)}$ is the union of the irreducible components of $X$, the assertion about the residue fields in \labelcref{lemitem:surgery-finite-normalisation} follows at once, and assertion \labelcref{lemitem:surgery-tameness} is also an immediate consequence. 

For assertion \labelcref{lemitem:surgery-fs-and-smooth}, we note first that for faces $F_1$ and $F_2$ in the image of $W \to \Spec(P)$, the intersection $F_{12} = F_1 \cap F_2$ is a face, and naturally 
\[
Z_{F_1} \cap Z_{F_2} = Z_{F_1} \times_W Z_{F_2} = Z_{F_{12}}.
\]
Consequently, we have the equality $X_{F_{12}} = X_{F_1} \cap X_{F_2}$. It follows that $X^{(a)}$ is also a disjoint union of normal subschemes of $X$ of the form $X_F$ for some faces $F$ of $P$. This makes $X^{(a)}_\star$ a disjoint union of subschemes $X_{F,\star}$ which are strict \'etale over $Z_{F,\star}$ and therefore fs and log smooth over $\Spec(k)$ with trivial log structure. 
 
The map $X^{(a)}\to X^{(a)}_\star$ in \labelcref{lemitem:surgery-relative-strip-off-log} is locally on $X_F$ given by  the map $X_F\to X_{F,\star}$ induced by the obvious map 
\[
(F \to k[F]) \la (P \to k[F])
\]
of prelog rings. In order to construct the map globally we need to establish its uniqueness locally. 
Since, by definition, $\cM_{X^{(a)}_\star}$ is a subsheaf of $\cO_{X^{(a)}}$, it suffices to show that $\alpha\colon \cM_{X^{(a)}}\to\cO_{X^{(a)}}$ maps the preimage of $\cM_{X^{(a)}_\star}\subseteq \cO_{X^{(a)}}$ isomorphically onto $\cM_{X^{(a)}_\star}$. This is an \'etale  local question to be checked in a geometric point $\bar{x}$ of $X_F$.  An element $(q,u)$ of the stalk 
\[
\cM_{X^{(a)},\bar x} = P\oplus_{\alpha_{\bar x}^{-1}(\cO_{X,(\bar x)}^\times)} \cO_{X,(\bar x)}^\times
\]
maps to a nonzero element of $\cO_{X,(\bar x)}$ if and only if $q\in F$. Thus, the preimage in question is $F\oplus_{\alpha_{\bar x}^{-1}(\cO_{X,(\bar x)}^\times)} \cO_{X,(\bar x)}^\times$, which maps isomorphically onto the stalk $\cM_{X^{(a)}_\star, \bar x}$. This shows \labelcref{lemitem:surgery-relative-strip-off-log} locally, and thus constructs the map in question globally together with the claimed uniqueness.
\end{proof}

\begin{rmk}
    The proof of \labelcref{lemitem:surgery-relative-strip-off-log} used the special local form of $X^{(a)}$, and the map $\varepsilon$ might not exist or be non-unique in simple situations. For an example exhibiting non-uniqueness, consider the diagonal in $\bA^2 = \Spec(\bN^2\to k[\bN^2])$, i.e.\ $X=\Spec(\bN^2\to k[x])$ where $\bN^2\to k[x]$ sends both generators $e_0,e_1$ to $x$. In this case, $U=D(x)$ and $X_\star = \Spec(\bN\to k[x])$ where $\bN\to k[x]$ sends the generator $e$ to $x$. There are two different maps $\varepsilon_i\colon X \to X_\star$, for $i=0,1$, induced by the identity on $k[x]$ and the two inclusions $\bN\to\bN^2$ sending $e\mapsto e_i$.
\end{rmk}

Thanks to \cref{lem:surgery-properties}~\ref{lemitem:surgery-finite-normalisation} and \ref{lemitem:surgery-tameness} and \cref{cor:descent-beyond-fs}, the van Kampen formula of \cref{sec:descent for Gal cats} yields (up to a choice of base points and paths) an isomorphism
\[ 
    \pitame(X/k) \simeq \pitame(\Desc(X^{(0)}/X), \FEtt_{-/k}).
\]
The next lemma allows us to strip off part of the log structure on the connected components of $X^{(a)}$ appearing implicitly in the above formula. It is a ``partially compactified'' version of \cref{cor:torus-fibration-tame}.

\begin{lem} \label{lem:right-exact-Xa}
    Let $Y$ be a connected component of $X^{(a)}$ and let $Y_\star$ be the corresponding component of $X^{(a)}_\star$ (the underlying schemes of $X^{(a)}$ and $X^{(a)}_\star$ are the same). Let $\bar\eta$ be a geometric generic point of $Y$. Then we have an exact sequence
    \[
        \begin{tikzcd}
            \pi_1(\overline{\cM}_{Y,\bar\eta}) \ar[r] & \pitame(Y/k,\bar\eta)\ar[r] & \pitame(Y_\star/k,\bar\eta)\ar[r] & 1.
        \end{tikzcd}
    \]
    The group on the left is isomorphic to $\widehat{\mathbb{Z}}'(1)^r$ for some $r\geq 0$. 
\end{lem}

\begin{proof}
Let $U=U^{(a)}\cap Y$ be the locus where $\overline{\cM}_Y$ is locally constant, let $j\colon U\to Y$ be the inclusion, and let $U_\star\subseteq Y_\star$ be the corresponding open. Then $U_\star=\underline{U}$ is the locus where the log structure on $Y_\star$ is trivial, and since $Y_\star$ is log regular by \cref{lem:surgery-properties}~\labelcref{lemitem:surgery-fs-and-smooth}, the open $U_\ast$ is connected and the map $\pitame(U_\star/k)\to \pitame(Y_\star/k)$ is an isomorphism by log purity as recalled later in \cref{thm:GR-log-Abhyankar}. We have a commutative diagram of profinite groups
\[
    \begin{tikzcd}
            \pi_1(\overline{\cM}_{U,\bar\eta}) \ar[r] \ar[d,equal] & \pitame(U/k,\bar\eta)\ar[r] \ar[d] & \pitame(U_\star/k,\bar\eta)\ar[r] \ar[d,"\simeq\text{ by purity}"] & 1 \\
            \pi_1(\overline{\cM}_{Y,\bar\eta}) \ar[r] & \pitame(Y/k,\bar\eta)\ar[r] & \pitame(Y_\star/k,\bar\eta)\ar[r] & 1
    \end{tikzcd}
\]
in which the top row is exact by \cref{cor:torus-fibration-tame}. We claim that the middle vertical arrow $\pitame(U/k)\to\pitame(Y/k)$ is surjective. If this is the case, then an easy diagram chase shows that the bottom row is exact as well, completing the proof.

The assertion that $\pitame(U/k)\to\pitame(Y/k)$ is surjective is equivalent to the full faithfulness of $\FEtt_{Y/k}\to\FEtt_{U/k}$. Since both source and target are full subcategories of $\FEt$, we can ignore the tameness condition and show the full faithfulness of $\FEt_Y\to\FEt_U$. This translates to: if $Y'\in \FEt_Y$ is connected, then $U' = Y'\times_Y U$ is connected. In order to show this, it is enough to show that $U'$ is dense in $Y'$ and that $Y'$ is geometrically unibranch. We can check these properties by an étale local calculation. Locally, $X$ has the form $\Spec(P\to k[P]/(M_+))$ and $Y$ is a connected component of $\Spec(P\to k[F])$ for a face $F \subseteq P$ with $F \cap M = 0$, see the proof of \cref{lem:surgery-properties}. Then $Y'$ locally is a connected component of  
\[
    \Spec(Q \to k[F]\otimes_{k[P]} k[Q])
\]
for some Kummer \'etale map $P \to Q$ (\cref{lem:fKet-local-structure}). Let $G$ be the face of $Q$ corresponding to $F$ under the bijection $\Spec(Q) \to \Spec(P)$. Then the reduction of $Y'$ is locally of the form $\Spec(Q \to k[G])$. Now $F \to G$ is Kummer \'etale, hence $G$ is also fs and smooth over $0$.
Thus, the reduction of $Y'$ is normal, and hence $Y'$ is geometrically unibranch. The same calculation shows that $U'$ is dense as its underlying topological space locally equals that of $\Spec(k[G^{\rm gp}])$ in the above presentation.
\end{proof}

\begin{cor} \label{cor:surgery-fp}
    In the situation of \cref{setup:surgery}, suppose that $k$ is algebraically closed and that the smooth scheme $U^{(0)}$ (with trivial log structure) admits a projective toroidal  compactification. Then $\pitame(X/k)$ is finitely presented. 
\end{cor}

\begin{proof}
By the van Kampen formula of \S\ref{sec:descent for Gal cats}, since each $X^{(a)}$ (for $a\leq 2$) has finitely many connected components, it suffices to show that for every connected component $Y$ of some $X^{(a)}$, the group $\pitame(Y/k)$ is finitely generated (true by our general result \cref{thm:pi1tame log is fg}), and finitely presented if $a=0$. By \cref{lem:right-exact-Xa}, this group sits inside an exact sequence 
\[
    \begin{tikzcd}
        \pi_1(\overline{\cM}_{Y,\bar\eta}) \ar[r] & \pitame(Y/k,\bar\eta)\ar[r] & \pitame(Y_\star/k,\bar\eta)\ar[r] & 1.
    \end{tikzcd}
\]
Arguing as in the proof of \cref{cor:pi1tame locally constant log is fg}, we see that $\pitame(Y/k)$ is finitely presented if and only if $\pitame(Y_\star/k)$ is. On the other hand, $U=U^{(0)}\cap Y_\star$ (which is a connected component of $U^{(0)}$) is the locus where the log structure on $Y_\star$ is trivial, and hence by log purity (\cref{thm:GR-log-Abhyankar}) we have $\pitame(Y_\star/k)\simeq \pitame(U/k)$ (see the proof of \cref{lem:right-exact-Xa}). By our assumption, $U$ has a projective toroidal compactification $\overline{U}$. Applying a suitable log blowup \cite{Niziol}, we may assume that $\overline{U}$ is an snc compactification, in which case the assertion follows from \cref{thm:naked pi1tame fp}.
\end{proof}

\begin{ex}[Semistable curve, continued] \label{ex:semistable-log-curve-2}
    Consider the situation in \cref{ex:semistable-log-curve}, and assume that $k$ is algebraically closed. Let $Y_i = \Spec(P \to k[x_i])$, for $i = 1,2$, denote the two connected components of $X^{(0)}$.
    Since $M=V$ is divisible and equal to the stalk of $\overline{\cM}$ at the generic points of $X^{(0)}$, for $a=0$ the group $\pi_1(\overline{\cM}_{Y_i,\bar\eta})$ in \cref{lem:right-exact-Xa} is trivial.  This identifies both $\pitame(Y_1/k)$ and $\pitame(Y_2/k)$ canonically with $\pitame(\Spec(\bN \to k[x])/k) = \widehat{\mathbb{Z}}'(1)$. As both maps $Y_i \to X$ are monomorphisms, we can by \cite[Proposition~4.3]{Stix2006:vanKampen} simplify the category of descent data to `ordered descent data' where we fix an ordering of the connected components of $X^{(0)}$ and only consider those components of $X^{(a)}$ arising from fibre products for ordered $(a+1)$-many factors. Here this means we are left with just $Y_{12} = Y_1 \times_X Y_2 = \Spec(P \to k)$. \cref{prop:FKEt-pi1-local} yields the formula $\pitame(Y_{12}/k) = \pi_1(P) = \widehat{\mathbb{Z}}'(1)$.   
    By van Kampen, we have a pushout diagram
    \[
        \begin{tikzcd}
            \pitame(Y_{12}/k) \ar[r] \ar[d] & \pitame(Y_1/k) \ar[d] \\
            \pitame(Y_2/k) \ar[r] & \pitame(X/k)
        \end{tikzcd}
    \]
    which thanks to what we know equals
    \[
        \begin{tikzcd}
            \widehat{\mathbb{Z}}'(1) \ar[r] \ar[d] & \widehat{\mathbb{Z}}'(1) \ar[d] \\
            \widehat{\mathbb{Z}}'(1) \ar[r] & \pitame(X/k).
        \end{tikzcd}
    \]
    The top and left maps are isomorphisms (a special case of \cref{lem:pi1XF} below), and we conclude that 
    \[
        \pitame(X/k) \simeq \widehat{\mathbb{Z}}'(1).  
    \]
\end{ex}

The following lemma (whose inclusion was prompted by a question of Paul Alexander Helminck, see \cref{ex:helminck}) will allow us to treat more general examples of the toric kind. 

\begin{lem} \label{lem:pi1XF}  
     In the Zariski local situation of \cref{setup:surgery}, assume $k$ is algebraically closed and $P$ is sharp. Let $F \subseteq P$ be a face with $F \cap M = 0$, and denote 
     \[
     Z_F = \Spec(P\to k[F])
     \]
     as in \cref{lem:strata-closures}. Then $Z_F$ is connected and the map $Z_0 = \Spec(P\to k) \inj Z_F$, where the subscript $0$ refers to the face $0\subseteq P$, induces isomorphisms
     \[
        \pi_1(P) = \pi_1(Z_0) = \pitame(Z_0/k)\la \pitame(Z_F/k)
    \]
    (since these groups are abelian, we can omit base points). 
\end{lem}

\begin{proof}
Since $P$ is sharp, also all its faces $F \subseteq P$ are sharp, and $F^\gp$ does not contain torsion. Since $\Spec(k[F^\gp])$ is dense in $\underline{Z}_F$, the log scheme $Z_F$ is connected. 

Let us now show that $\FEtt_{Z_F/k}\to \FEtt_{Z_0/k} = \FEt_{Z_0}$ is essentially surjective. By \cref{prop:FKEt-pi1-local}, every connected cover of $Z_0$ is of the form $\Spec(Q\to k[Q]\otimes_{k[P]} k)$ for some Kummer \'etale map $P\to Q$. This map is the base change to $Z_0$ of the Kummer \'etale map
\[
    Y_F(Q) = \Spec(Q \to k[Q] \otimes_{k[P]} k[F]) \la Z_F.
\]
This map is Galois as an object of $\FEt_{Z_F}$ with Galois group $\Hom(Q^\gp/P^\gp,\mu_\infty(k))$ 
of degree prime to $p$, and hence is tame relative to $k$. 

To show full faithfulness, we need to show that if $Y \to Z_F$ is a connected cover which is tame relative to $k$, then $Y_0 = Y \times_{Z_F} Z_0$ is connected. To this end, consider the map 
\[
    (Y_{\rm red})_\star \la Z_{F,\star} = \Spec(F \to k[F]).
\]
By \cref{lem:fKet-local-structure}, $Y$ is \emph{locally on $Z_F$} isomorphic to the disjoint union of some $Z_F(Q)$'s. Since $(Z_F(Q)_{\rm red})_\star \simeq \Spec(G\to k[G])$ where $G \subseteq Q$ is the saturation of $F$ in $Q$ (or the unique face above $F$), see the computation in the proof of \cref{lem:right-exact-Xa}, it follows that $(Y_{\rm red})_\star$ is locally on $Z_{F,\star}$ the disjoint union of $\Spec(G \to k[G])$ for some Kummer \'etale maps $F \to G$. We deduce from this analysis that $(Y_{\rm red})_\star \to Z_{F,\star}$ is Kummer \'etale and tame. 
On the other hand, by purity as in the proof of \cref{lem:right-exact-Xa} we compute 
\[
\pitame(Z_{F,\star}/k) = \pitame(\Spec(k[F^{\rm gp}])/k) = \pi_1(F).
\]
It follows that 
$(Y_{\rm red})_\star \simeq \Spec(G\to k[G])$
for some Kummer \'etale map $F\to G$. The fibre of $(Y_{\rm red})_\star$ above $\underline{Z}_0=\Spec(k)$ consists of a single (possibly non-reduced) point, and hence is connected. However, the log schemes $(Y_{\rm red})_\star$ and $Y$ share their underlying topological space, and we conclude that $Y_0$ is connected.
\end{proof}

\begin{cor} \label{cor:Paul}
    In the Zariski local situation of \cref{setup:surgery}, where we consider the log scheme  
    \[
         W = \Spec(P\to k[P]\otimes_{k[M]} k) = \Spec\big(P\to k[P]/(\theta(M_+))\big),
    \]  
    assume $k$ is algebraically closed and $P$ is sharp. Then $W$ is connected and the inclusion $Z_0 \inj W$ (notation as in \cref{lem:strata-closures}) induces an isomorphism
    \[
        \pi_1(P) = \pi_1(Z_0) = \pitame(Z_0/k) \isomto \pitame(W/k). 
    \]   
\end{cor}

\begin{proof}
Since $P$ is sharp, also all its faces $F \subseteq P$ are sharp. Therefore the closures of log strata $Z_F \inj W$ are connected by \cref{lem:pi1XF}. Since all $Z_F$ contain the minimal log stratum $Z_0 = \Spec(P \to k)$, the scheme $W$ is connected. 

We apply the van Kampen formula with ordered descent data given by the  connected components $Z_F = \Spec(P \to k[F])$ of $W^{(0)}$ (with $F \subseteq P$ maximal faces in the image of $W \to \Spec(P)$, see the proof of \cref{lem:surgery-properties}). We deduce from \cref{lem:pi1XF} that all connected components of $W^{(a)}$, for all $a\geq 0$, have fundamental groups canonically isomorphic to $\pi_1(P)$.  Since all fibre products of ordered factors are again of the form $Z_F$ and so connected, the underlying $2$-complex for ordered descent data is the $2$-truncation of a full simplex. Moreover, all constituents have tame fundamental group canonically isomorphic to $\pi_1(P) = \pitame(Z_0/k)$ which is abelian (thus base points do not matter), so we can choose the paths so that the annoying datum $\alpha_{ijk}^{(f)}$ in \cite[Corollary~5.3]{Stix2006:vanKampen} is the identity in all cases. Thus, the van Kampen formula simplifies to just~$\pi_1(P)$.
\end{proof}

